\subsection{The solution of equations in a formal power series ring}

Lemma 2. --- Let \((g_i)_{i \in I}\) be a family of elements of order \(\geqslant 2\) in A{[}{[}I{]}{]}. When I is infinite, assume that \(g_i\) tends to 0 along the filter of complements of finite subsets of I. There exists one and only one automorphism T of the topological A-algebra A{[}{[}I{]}{]} such that T(X \(_i\) ) = X \(_i\) + g \(_i\) for all i ∈ I. Further,

\[
\omega (\mathrm{T} (u) - u) \geqslant \omega (u) + 1\tag{23}
\]

for each \(u \in \mathbf{A}[[\mathrm{I}]]\) .

The series \(f_{i} = \mathbf{X}_{i} + g_{i}\) has no constant term and when I is infinite, \(f_{i}\) tends to 0 along the filter of complements of finite subsets of I. Therefore (IV, p.~28, Prop. 4) there exists precisely one continuous endomorphism T of the A-algebra A{[}{[}I{]}{]} such that \(\mathrm{T}(\mathbf{X}_i) = f_i\) for all \(i \in I\) . For each \(\nu \in \mathbf{N}^{(I)}\) we put

\[
v _ {\nu} = \mathrm{T} \left(\mathbf {X} ^ {\nu}\right) - \mathbf {X} ^ {\nu} = \prod_ {i \in I} \left(\mathbf {X} _ {i} + g _ {i}\right) ^ {\nu (i)} - \prod_ {i \in I} \mathbf {X} _ {i} ^ {\nu (i)};
\]

the relations \(\omega(g_i) \geqslant 2\) imply \(\omega(v_\nu) \geqslant |\nu| + 1\) , and the relation (23) follows at once from this.

Let us show that T is injective. Given \(u \in A[[I]]\) such that \(T(u) = 0\) , by (23) we have \(\omega(u) \geqslant \omega(u) + 1\) , which is impossible if \(u \neq 0\) because \(\omega(u)\) would then be a positive integer.

For every formal series \(v\) in \(A[[I]]\) we denote by \(H_n(v)\) its homogeneous component of degree \(n\) . Let us put \(S_0(v) = H_0(v)\) and define the continuous mappings \(S_n: A[[I]] \to A[[I]]\) by the recursion equations

\[
\mathrm{S} _ {n} (v) = \mathrm{H} _ {n} \left(v - \mathrm{T} \left(\sum_ {k = 0} ^ {n - 1} \mathrm{S} _ {k} (v)\right)\right) \quad \text { for } \quad n \geqslant 1.\tag{24}
\]

Put \(S(v) = \sum_{n \geqslant 0} S_n(v)\) ; if \(\nu \in \mathbf{N}^{(I)}\) and \(n = |\nu|\) , then the coefficient \(S^{\nu}(v)\) of

\(X^{\nu}\) in \(S(v)\) is equal to that of \(X^{\nu}\) in \(S_{n}(v)\) ; since \(S_{n}\) is a continuous mapping, the mapping \(S^{\nu}: A[[I]] \to A\) is continuous. Hence, by the definition of the product topology on \(A[[I]] = A^{N^{(I)}}\) , the mapping \(S: A[[I]] \to A[[I]]\) is continuous.

We shall prove the relation \(\mathrm{T}(\mathrm{S}(v)) = v\) for all \(v \in A[[I]]\) which will complete the proof of the lemma. Let \(v \in A[[I]]\) , \(u_n = S_n(v)\) and \(u = S(v)\) . Let n be a positive integer such that

\[
\omega (v - \mathrm{T} (u)) \geqslant n.\tag{$(25)_{n}$}
\]

We have \(\omega(u - (u_{0} + \cdots + u_{n-1})) \geqslant n\) , whence

\[
\omega (\mathrm{T} (u) - \mathrm{T} (u _ {0} + \dots + u _ {n - 1}) - u _ {n}) \geqslant n + 1\tag{26}
\]

by (23). Now the recursion equation (24) shows that

\[
u _ {n} = \mathrm{H} _ {n} (v - \mathrm{T} (u _ {0} + \dots + u _ {n - 1}))\tag{27}
\]

By (26) the formal power series \(v - \mathrm{T}(u)\) and \(v - \mathrm{T}(u_{0} + \cdots + u_{n-1}) - u_{n}\) have the same homogeneous component of degree n, and this component is zero, by (27). We thus have \(\omega(v - \mathrm{T}(u)) \geqslant n + 1\) , that is \((25)_{n}\) implies \((25)_{n+1}\) . Since \((25)_{0}\) clearly holds, we thus have \(\omega(v - \mathrm{T}(u)) \geqslant n\) for every integer \(n \geqslant 0\) , whence \(v = \mathrm{T}(u) = \mathrm{T}(\mathrm{S}(v))\) , as was to be proved.

For the rest of this No.~we shall, for any set I, denote by A \{I\} the set of families \((f_i)_{i\in I}\) satisfying the following conditions:

\begin{enumerate}
\def\labelenumi{(\roman{enumi})}
\item
  for each \(i \in I\) , \(f_i\) is an element of A{[}{[}I{]}{]} without constant term;
\item
  if I is infinite, \(f_{i}\) tends to 0 along the filter of complements of finite subsets of I.
\end{enumerate}

The set A\{I\} is a monoid for the composition law \((\mathbf{f},\mathbf{g})\mapsto\mathbf{f}\circ\mathbf{g}\) , with \(\{X_{i}\}_{i\in I}\) as unit element. The set of invertible elements of A\{I\} is thus a group.

On the other hand, let E be the monoid of all continuous unital endomorphisms of the A-algebra \(\mathbf{A}[[\mathrm{I}]]\) leaving the ideal of all formal power series without constant term invariant. If \(\mathbf{f} \in \mathbf{A}\{\mathbf{I}\}\) and \(g \in \mathbf{A}[[\mathrm{I}]]\) , then the element \(g(\mathbf{f})\) is defined. For fixed \(\mathbf{f}\) , the mapping \(g \mapsto g(\mathbf{f})\) of \(\mathbf{A}[[\mathrm{I}]]\) into itself is an element \(\mathbf{W}_{\mathbf{f}}\) of E. If \(\mathbf{f}_1, \mathbf{f}_2 \in \mathbf{A}\{\mathbf{I}\}\) and \(g \in \mathbf{A}[[\mathrm{I}]]\) , we have, by formula (5) (IV, p.~30)

\[
\mathbf {W} _ {\mathbf {f} _ {1} \circ \mathbf {f} _ {2}} (g) = g (\mathbf {f} _ {1} \circ \mathbf {f} _ {2}) = g (\mathbf {f} _ {1}) \circ \mathbf {f} _ {2} = \mathbf {W} _ {\mathbf {f} _ {2}} (\mathbf {W} _ {\mathbf {f} _ {1}} (g))
\]

hence \(\mathbf{f} \mapsto \mathbf{W}_{\mathbf{f}}\) is a homomorphism of the monoid opposite to A \{I\} into E. By Prop. 4 (IV, p.~28) this homomorphism is bijective.

Let \(\mathbf{f} = (f_i)_{i\in I}\in \mathbf{A}\{\mathbf{I}\}\) and let \(\sum_{j\in I}\alpha_{ij}\mathbf{X}_j\) be the homogeneous component of

degree 1 of \(f_{i}\) . For any fixed \(j\) in I we have \(\alpha_{ij} = 0\) except for a finite number of suffixes \(i\) , by hypothesis (ii) above. If \((\lambda_i) \in \mathbf{A}^{(I)}\) , we thus have \(\left(\sum_{j \in I} \alpha_{ij} \lambda_j\right) \in \mathbf{A}^{(I)}\) .

We denote by \(T_{f}\) the A-linear mapping \(^{1}\)

\[
\left(\lambda_ {i}\right) \mapsto \left(\sum_ {j \in I} \alpha_ {i j} \lambda_ {j}\right)
\]

of \(\mathbf{A}^{(I)}\) into \(\mathbf{A}^{(I)}\) . If \(\mathbf{g} \in \mathbf{A}\{I\}\) , it is easily verified that

\[
\mathrm{T} _ {\mathbf {f} \circ \mathbf {g}} = \mathrm{T} _ {\mathbf {f}} \circ \mathrm{T} _ {\mathbf {g}}.\tag{28}
\]

PROPOSITION 10. --- Let \(\mathbf{f} \in \mathbf{A}\{\mathbf{I}\}\) ; then the following conditions are equivalent: (i) \(\mathbf{f}\) is invertible in \(\mathbf{A}\{\mathbf{I}\}\) for the law \(\circ\) ;

\begin{enumerate}
\def\labelenumi{(\roman{enumi})}
\setcounter{enumi}{1}
\tightlist
\item
  \(\mathbf{T}_{\mathbf{f}}\) is invertible in the ring \(\operatorname{End}(\mathbf{A}^{(I)})\) .
\end{enumerate}

The implication (i) \(\Rightarrow\) (ii) is immediate from (28). Suppose now that \(T_{f}\) is invertible in \(\operatorname{End}(A^{(I)})\) . There exists \(\mathbf{g} = (g_i)_{i \in I} \in A\{I\}\) such that each \(g_i\) is homogeneous of degree 1 and \(T_{g} \circ T_{f}\) is the identity mapping of \(A^{(I)}\) . Write \(\mathbf{h} = \mathbf{g} \circ \mathbf{f}\) ; then (28) shows that \(T_{h}\) is the identity mapping of \(A^{(I)}\) , which is equivalent to the assertion \(\omega(h_i - X_i) \geqslant 2\) . By Lemma 2 of IV, p.~35 h is therefore invertible in A \{I\}. It is clear that \(\mathbf{g}\) is invertible in A \{I\}, hence \(\mathbf{f}\) is invertible in A \{I\}.

COROLLARY. --- Let \(f_{i}(Y_{1}, Y_{2}, \ldots, Y_{q}, X_{1}, X_{2}, \ldots, X_{p}) (1 \leqslant i \leqslant q)\) be q formal power series without constant term in A{[}{[} \(Y_{1}, \ldots, Y_{q}, X_{1}, \ldots, X_{p}\) {]}{]}. If the constant term of the formal power series D = det \(\left(\frac{\partial f_{i}}{\partial Y_{j}}\right)\) is invertible in A, then there exists precisely one system of q formal power series \(u_{1}(X_{1}, \ldots, X_{p}), \ldots, u_{q}(X_{1}, \ldots, X_{p})\) such that

\[
f _ {i} (u _ {1}, \dots , u _ {q}, \mathrm{X} _ {1}, \dots , \mathrm{X} _ {p}) = 0 \quad (1 \leqslant i \leqslant q).\tag{29}
\]

Put \(f_{q+1} = \mathbf{X}_1, \ldots, f_{q+p} = \mathbf{X}_p\) , \(\mathbf{f} = (f_1, \ldots, f_{p+q})\) , then \(\det \mathbf{T}_{\mathbf{f}}\) is equal to the constant term of \(\mathbf{D}\) , hence invertible in \(\mathbf{A}\) ; therefore \(\mathbf{T}_{\mathbf{f}}\) is invertible. By Prop. 10 there exist formal power series without constant term

\[
\boldsymbol {g} _ {1}, \dots , \boldsymbol {g} _ {q + p} \in \mathrm{A} \left[ \left[ \mathrm{Y} _ {1}, \dots , \mathrm{Y} _ {q}, \mathrm{X} _ {1}, \dots , \mathrm{X} _ {p} \right] \right]
\]

such that on writing

\[
\mathbf {g} = \left(g _ {1}, \dots , g _ {p + q}\right), \quad \mathbf {1} _ {p + q} = \left(\mathrm{Y} _ {1}, \dots , \mathrm{Y} _ {q}, \mathrm{X} _ {1}, \dots , \mathrm{X} _ {p}\right)
\]

we have \(\mathbf{f} \circ \mathbf{g} = \mathbf{g} \circ \mathbf{f} = \mathbf{1}_{p + q}\) . The relation \(\mathbf{f} \circ \mathbf{g} = \mathbf{1}_{p + q}\) in particular gives

\[
\boldsymbol {g} _ {q + 1} = \mathbf {X} _ {1}, \dots , \boldsymbol {g} _ {q + p} = \mathbf {X} _ {p}.
\]

Hence

\[
f _ {i} \left(g _ {1}, \dots , g _ {q}, \mathrm{X} _ {1}, \dots , \mathrm{X} _ {p}\right) = \mathrm{Y} _ {i} \quad (1 \leqslant i \leqslant q).\tag{30}
\]

Now put

\[
u _ {i} \left(\mathrm{X} _ {1}, \dots , \mathrm{X} _ {p}\right) = g _ {i} (0, \dots , 0, \mathrm{X} _ {1}, \dots , \mathrm{X} _ {p}) \quad (1 \leqslant i \leqslant q);\tag{31}
\]

substituting 0 for each \(\mathbf{Y}_i\) in (30) we obtain the desired relation (29).

Conversely, suppose that the formal power series \(u_1, \ldots, u_q\) in the ring \(A[[X_1, \ldots, X_p]]\) satisfy the relation (29). The relation \(\mathbf{g} \circ \mathbf{f} = 1_{p + q}\) implies

\[
g _ {i} (f _ {1}, \dots , f _ {q}, \mathrm{X} _ {1}, \dots , \mathrm{X} _ {p}) = \mathrm{Y} _ {i} \quad (1 \leqslant i \leqslant q);\tag{32}
\]

and substituting \(u_{i}\) for \(Y_{i}\) for \(1 \leqslant i \leqslant q\) in (32), we obtain (31), whence the uniqueness of the solution of the system (29).

